\subsection{Performance with Different Backbone LLMs}

\begin{table}[t]
% \vspace{-10pt}
\caption{A comparison of \methodname{} with different backbone LLMs against three baselines on the MatplotBench benchmark. All results are presented in percent (\%).}
\centering
\footnotesize 
\setlength{\tabcolsep}{6.5pt} 
\begin{adjustbox}{max width=\linewidth}
\begin{tabular}{@{}l|ccc |ccc |ccc@{}}
\toprule
\textbf{Base LLMs}& \multicolumn{3}{c|}{\textbf{Gemini-2.5-Pro}} & \multicolumn{3}{c|}{\textbf{Gemini-2.5-Flash}} & \multicolumn{3}{c}{\textbf{Claude-4-Sonnet}} \\
\midrule
\textbf{Method} & EPR $\uparrow$ & VSR $\uparrow$ & OS $\uparrow$ & EPR $\uparrow$ & VSR $\uparrow$ & OS $\uparrow$ & EPR $\uparrow$ & VSR $\uparrow$ & OS $\uparrow$ \\
\midrule
MatplotAgent & 92.0 & 55.4 & 51.0 & 99.0 & 46.4 & 45.9 & 93.0 & 58.8 & 54.7 \\
VisPath & 73.0 & 60.5 & 44.2 & 95.0 & 45.8 & 43.5 & 57.0 & \best{77.5} & 44.2 \\
CoML4VIS & 99.0 & 63.2 & 62.6 & 99.0 & 57.8 & 57.2 & 99.0 & 65.9 & 65.2 \\ \midrule
\methodname{} (Ours) & \best{99.0} & \best{80.3} & \best{79.5} & \best{99.0} & \best{78.5} & \best{77.7} & \best{98.0} & 76.7 & \best{75.2} \\
\bottomrule
\end{tabular}
\end{adjustbox}
\label{tab:backbone-comparison}
% \vspace{-5pt}
\end{table}

To assess the generality of \methodname{} across diverse LLM backbones, we evaluate its performance when substituting the primary \io{gemini-2.5-pro} with alternative strong capability LLMs: \io{gemini-2.5-flash} and \io{claude-4-sonnet}. 
This experiment isolates the impact of the backbone LLM on visualization generation, holding constant the multi-agent architecture. We focus on the MatplotBench, as it emphasizes robust handling of numerical data, pattern recognition, and code synthesis under ambiguous queries—tasks that stress the backbone's reasoning and code generation capabilities.


We select these backbones for their complementary strengths: \io{gemini-2.5-flash} prioritizes efficiency and low-latency inference, making it suitable for real-time applications, while \io{claude-4-sonnet} excels in language understanding and multi-step reasoning, potentially enhancing agent collaboration in complex scenarios. All models are configured with identical hyperparameters. Table~\ref{tab:backbone-comparison} presents the results. \methodname{} with \io{gemini-2.5-flash} achieves an OS of 77.7\%, showcasing efficient handling of real-time scenarios with minimal degradation (1.8\% relative to \io{gemini-2.5-pro}), attributable to streamlined agent interactions that leverage metadata over raw data ingestion. \io{claude-4-sonnet}, conversely, attains an OS of 75.2\%, a 4.3\% drop from \io{gemini-2.5-pro} , likely stemming from its enhanced semantic parsing but reduced robustness in code execution under high-context loads. These outcomes highlight \methodname{}'s backbone-agnostic design, amplifying each LLM's inherent strengths while mitigating weaknesses through collaborative workflows. 

We compare \methodname{} against each other using the three backbone LLMs as described above. Across the board, \methodname{} outperforms baselines significantly, with the best-performing variant, \methodname{} with \io{gemini-2.5-pro} , achieving 79.5\% OS. \textit{MatplotAgent}, \textit{VisPath}, and \textit{CoML4VIS} struggle to exceed 65.2\% OS in any setting, highlighting the challenges of visualization tasks without multi-agent refinement. We also observe that \methodname{} trends similarly across different backbones, with EPR and VSR remaining consistently high (98.0--99.0\% and 76.7--80.3\%). 


LLMs tend to generate simpler visualizations. Baseline-generated code tends to produce fewer refinements than \methodname{}. As shown in Table~\ref{tab:backbone-comparison}, compared to \methodname{}, baselines like \textit{MatplotAgent} achieve lower VSR (46.4--58.8\%), and rarely handle complex multi-faceted queries.

